% --- Содержимое файла materials/4_4.tex ---

\subsection{Элемент наилучшего приближения. Теорема о проекции}

\begin{definition}[Элемент наилучшего приближения]
	Пусть $E$ --- линейное нормированное пространство, $L \subset E$ --- подпространство, $h \in E$. \textbf{Элементом наилучшего приближения} для $h$ называется вектор $x \in L$ такой, что расстояние от $h$ до $x$ минимально среди всех векторов $L$:
	\[ \|h - x\| = \rho(h, L) = \inf \limits_{y \in L} \|h - y\|. \]
\end{definition}

\begin{remark}
	В общем случае элемент наилучшего приближения может не существовать (если подпространство не замкнуто или пространство не рефлексивно) или быть не единственным (если норма "угловатая").
\end{remark}

\begin{proposition}[Существование и единственность в Гильбертовом пространстве]\label{bestelt}
	Пусть $H$ --- гильбертово пространство, $M \subset H$ --- \textbf{замкнутое} подпространство. Тогда для любого $h \in H$ существует единственный элемент наилучшего приближения $x \in M$.
\end{proposition}

\begin{proof}
	Зафиксируем $h \in H$. Обозначим $d := \rho(h, M)$.
	
	\textbf{1. Существование.}
	\begin{itemize}
		\item Если $d = 0$, то в силу замкнутости $M$, $h \in M$, и искомый элемент — сам $h$.
		\item Пусть $d > 0$. По определению инфимума выберем минимизирующую последовательность $\{x_n\} \subset M$ такую, что $\|h - x_n\| \to d$.
		\item Проверим фундаментальность $\{x_n\}$. Применим равенство параллелограмма к векторам $(h - x_n)$ и $(h - x_m)$:
		\[ \|(h - x_n) + (h - x_m)\|^2 + \|(h - x_n) - (h - x_m)\|^2 = 2\|h - x_n\|^2 + 2\|h - x_m\|^2 \]
		Преобразуем:
		\[ \|2h - (x_n + x_m)\|^2 + \|x_m - x_n\|^2 = 2\|h - x_n\|^2 + 2\|h - x_m\|^2 \]
		\[ 4\left\|h - \frac{x_n + x_m}{2}\right\|^2 + \|x_n - x_m\|^2 = 2\|h - x_n\|^2 + 2\|h - x_m\|^2 \]
		\item Заметим, что $\frac{x_n + x_m}{2} \in M$, следовательно $\left\|h - \frac{x_n + x_m}{2}\right\| \geqslant d$.
		\item Переходим к пределу. Правая часть стремится к $2d^2 + 2d^2 = 4d^2$. Левая часть $\geqslant 4d^2 + \|x_n - x_m\|^2$. Отсюда следует, что $\|x_n - x_m\| \to 0$.
		\item Так как $H$ полно и $M$ замкнуто, $x_n \to x \in M$. В силу непрерывности нормы $\|h - x\| = d$.
	\end{itemize}
	
	\textbf{2. Единственность.}
	\begin{itemize}
		\item Пусть существуют два элемента $x, y \in M$, реализующих минимум: $\|h - x\| = d$ и $\|h - y\| = d$.
		\item Применим равенство параллелограмма к векторам $h-x$ и $h-y$:
		\[ \|(h-x) + (h-y)\|^2 + \|(h-x) - (h-y)\|^2 = 2\|h-x\|^2 + 2\|h-y\|^2 \]
		\[ 4\left\|h - \frac{x+y}{2}\right\|^2 + \|y-x\|^2 = 2d^2 + 2d^2 = 4d^2 \]
		\item Так как $\frac{x+y}{2} \in M$, то $\left\|h - \frac{x+y}{2}\right\| \geqslant d$, а значит первое слагаемое $\geqslant 4d^2$.
		\item Получаем неравенство $4d^2 + \|x-y\|^2 \leqslant 4d^2$, откуда $\|x - y\| = 0 \Rightarrow x = y$.
	\end{itemize}
\end{proof}

\begin{idea}[О наилучшем приближении]
	
	\textbf{Смысл утверждения:}
	
	Представьте, что $M$ — это бесконечная плоская поверхность (подпространство), а $h$ — точка, висящая над ней. Утверждение гарантирует, что на этой поверхности есть \textbf{ровно одна} точка $x$, ближайшая к $h$ (проекция).
	
	\begin{itemize}
		
		\item В «плохих» пространствах (как $L_\infty$ с квадратными шарами) ближайших точек может быть много (целая грань квадрата).
		
		\item В Гильбертовом пространстве шары «круглые» (строго выпуклые), поэтому касание происходит только в одной точке.
		
	\end{itemize}
	
	\textbf{Идея доказательства:}
	
	\begin{enumerate}
		
		\item \textbf{Существование:} Мы берем последовательность точек, которые подходят всё ближе к идеальному расстоянию $d$. Чтобы доказать, что они не «танцуют» на месте, а сходятся к цели, мы используем \textbf{равенство параллелограмма}. Оно переводит геометрическое сближение (расстояния) в алгебраический факт фундаментальности.
		
		\item \textbf{Единственность:} Если бы было две наилучшие точки $x$ и $y$, то середина отрезка между ними $\frac{x+y}{2}$ оказалась бы \textit{еще ближе} к цели (из-за выпуклости шара), что противоречит тому, что $x$ и $y$ — уже самые близкие.
		
	\end{enumerate}
	
\end{idea}

\begin{definition}[Аннулятор]
	Пусть $E$ --- евклидово пространство, $S \subset E$. \textbf{Аннулятором} (или ортогональным дополнением) множества $S$ называется множество:
	\[S^\perp := \{y \in E \mid \forall x \in S: (x, y) = 0\}\]
\end{definition}

\begin{remark}
	$S^\perp$ всегда является \textbf{замкнутым подпространством} в $E$ (даже если $S$ не замкнуто и не подпространство).
	Справедливо: $S^\perp = [S]^\perp = \overline{[S]}^\perp$.
\end{remark}

\begin{theorem}[Теорема Рисса о проекции]\label{thm4.3}
	Пусть $H$ --- гильбертово пространство, $M \subset H$ --- \textbf{замкнутое} подпространство. Тогда пространство разлагается в прямую ортогональную сумму:
	\[ H = M \oplus M^\perp \]
	Это означает, что любой вектор $h \in H$ единственным образом представим в виде $h = x + y$, где $x \in M, y \in M^\perp$.
\end{theorem}

\begin{idea}
	Теорема утверждает, что мы можем "уронить перпендикуляр" из любой точки на подпространство. Вектор $x$ — это проекция ("тень"), а вектор $y$ — перпендикуляр ("высота").
\end{idea}


\begin{proof}
	\textbf{1. Возможность разложения ($H = M + M^\perp$).}
	\begin{itemize}
		\item Зафиксируем $h \in H$. По предложению \ref{bestelt}, существует единственный элемент наилучшего приближения $x \in M$.
		\item Обозначим вектор невязки $y := h - x$. Покажем, что $y \in M^\perp$.
		\item Обозначим $d = \|y\| = \|h - x\|$. Для любого вектора $m \in M$ и любого числа $\alpha \in \mathbb{R}$ вектор $(x + \alpha m)$ лежит в $M$.
		\item Так как $x$ — ближайший элемент, расстояние до него минимально:
		\[ \|y\|^2 \leqslant \|h - (x + \alpha m)\|^2 = \|y - \alpha m\|^2 = (y - \alpha m, y - \alpha m) \]
		\[ \|y\|^2 \leqslant \|y\|^2 - 2\alpha(y, m) + \alpha^2\|m\|^2 \]
		\item Сокращаем $\|y\|^2$ и получаем: $2\alpha(y, m) \leqslant \alpha^2\|m\|^2$.
		\item Если бы $(y, m) \neq 0$, мы могли бы выбрать $\alpha$ очень малым и такого знака, что левая часть (линейная по $\alpha$) перевесила бы правую (квадратичную по $\alpha$), нарушив неравенство. (Формально: положим $\alpha = \text{sign}(y,m) \cdot \varepsilon$, при $\varepsilon \to 0$ получаем противоречие).
		\item Следовательно, $(y, m) = 0$ для всех $m \in M$. Значит $y \perp M$.
		\item Мы представили $h = x + y$, где $x \in M, y \in M^\perp$.
	\end{itemize}
	
	\textbf{2. Единственность (Прямая сумма).}
	\begin{itemize}
		\item Пусть $z \in M \cap M^\perp$.
		\item Так как $z \in M^\perp$ и $z \in M$, он должен быть ортогонален самому себе: $(z, z) = 0$.
		\item По аксиоме скалярного произведения $\|z\|^2 = 0 \Rightarrow z = 0$.
		\item Пересечение тривиально, значит сумма прямая.
	\end{itemize}
\end{proof}

\begin{remark}
	Это свойство (существование проекции) выделяет гильбертовы пространства среди всех банаховых. В общем случае в банаховом пространстве нельзя гарантировать существование такого "удобного" разложения.
\end{remark}

\begin{idea}[О теореме Рисса]
	
	\textbf{Смысл утверждения:}
	
	Любой вектор можно разложить на две независимые составляющие: «тень» на подпространстве ($x \in M$) и «высоту» над ним ($y \in M^\perp$). Это обобщение школьного разложения сил на нормальную и тангенциальную составляющие.
	
	\textbf{Идея доказательства (Вариационный метод):}
	
	Мы уже знаем, что $x$ — это самая близкая точка в $M$ к вектору $h$. Вектор ошибки (высота) — это $y = h - x$.
	
	Почему $y$ обязан быть перпендикулярен $M$?
	
	\begin{itemize}
		
		\item Представьте, что вектор $y$ наклонен к плоскости $M$ под углом $80^\circ$ (не перпендикулярен).
		
		\item Тогда мы могли бы чуть-чуть сместиться по плоскости в сторону этого наклона и уменьшить длину вектора $y$.
		
		\item Но $x$ — это уже \textbf{самая} близкая точка, уменьшать расстояние некуда.
		
		\item Значит, никакого наклона нет, угол ровно $90^\circ$.
		
	\end{itemize}
	
\end{idea}